% !TEX TS-program = pdflatex
% !TEX encoding = UTF-8 Unicode

% This file is a template using the "beamer" package to create slides for a talk or presentation
% - Giving a talk on some subject.
% - The talk is between 15min and 45min long.
% - Style is ornate.

% MODIFIED by Jonathan Kew, 2008-07-06
% The header comments and encoding in this file were modified for inclusion with TeXworks.
% The content is otherwise unchanged from the original distributed with the beamer package.

\documentclass{beamer}
\setbeamercovered{invisible}



% Copyright 2004 by Till Tantau <tantau@users.sourceforge.net>.
%
% In principle, this file can be redistributed and/or modified under
% the terms of the GNU Public License, version 2.
%
% However, this file is supposed to be a template to be modified
% for your own needs. For this reason, if you use this file as a
% template and not specifically distribute it as part of a another
% package/program, I grant the extra permission to freely copy and
% modify this file as you see fit and even to delete this copyright
% notice. 


\mode<presentation>
{
  \usetheme{Warsaw}
  % or ...

  % or whatever (possibly just delete it)
}


\usepackage[english]{babel}
% or whatever

\usepackage[utf8]{inputenc}
% or whatever

\usepackage{times}
\usepackage[T1]{fontenc}
% Or whatever. Note that the encoding and the font should match. If T1
% does not look nice, try deleting the line with the fontenc.

%%% MATH RELATED

\input{macros.tex}

\title{MATH 350-2 Advanced Calculus} 
\subtitle
{} % (optional)

\author[W.R. Casper] % (optional, use only with lots of authors)
{W.R. Casper}
% - Use the \inst{?} command only if the authors have different
%   affiliation.

\institute[California State University Fullerton] % (optional, but mostly needed)
{
  Department of Mathematics\\
  California State University Fullerton}
% - Use the \inst command only if there are several affiliations.
% - Keep it simple, no one is interested in your street address.

\subject{Talks}
% This is only inserted into the PDF information catalog. Can be left
% out. 



% If you have a file called "university-logo-filename.xxx", where xxx
% is a graphic format that can be processed by latex or pdflatex,
% resp., then you can add a logo as follows:

% \pgfdeclareimage[height=0.5cm]{university-logo}{university-logo-filename}
% \logo{\pgfuseimage{university-logo}}



% Delete this, if you do not want the table of contents to pop up at
% the beginning of each subsection:
\AtBeginSubsection[]
{
  \begin{frame}<beamer>{Outline}
    \tableofcontents[currentsection,currentsubsection]
  \end{frame}
}


% If you wish to uncover everything in a step-wise fashion, uncomment
% the following command: 

%\beamerdefaultoverlayspecification{<+->}


\begin{document}

\begin{frame}
  \titlepage
\end{frame}

\begin{frame}{Outline}
  \tableofcontents
  % You might wish to add the option [pausesections]
\end{frame}

% Since this a solution template for a generic talk, very little can
% be said about how it should be structured. However, the talk length
% of between 15min and 45min and the theme suggest that you stick to
% the following rules:  

% - Exactly two or three sections (other than the summary).
% - At *most* three subsections per section.
% - Talk about 30s to 2min per frame. So there should be between about
%   15 and 30 frames, all told.

\section{Real Analysis Lecture 24}

\subsection{Riemann-Stieltjes integrals}

\begin{frame}{Riemann sums}
A Riemann sum of $f(x)$ on $[a,b]$ with $n$ is
$$R(P,f,\{t_k\}) = \sum_{k=1}^n f(t_k)\Delta x_k$$
\pause
where $t_k$ and $\Delta x_k := x_k-x_{k-1}$ come from:
\begin{itemize}
\pause
\item a \textbf{partition} $P =\{x_0=a,x_1,\dots, x_n=b\}$
\pause
\item a set of $n$ \textbf{sample points} $t_k\in [x_{k-1},x_k]$ with $1\leq k\leq n$
\end{itemize}
\end{frame}

\begin{frame}{Riemann integrability}
\begin{defn}[Very informal first try]
The Riemann integral is the limit of \emph{all} Riemann sums as we increase the number of rectangles.
\end{defn}
\pause
\begin{quest}
What problems do you see with this definition?
\end{quest}
\end{frame}

\begin{frame}{Some issues}
Problems with definition:
\begin{itemize}
\item for $n$ rectangles, there are uncountably many partitions!
\item for a specific partition, there are uncountably many sample points!
\end{itemize}
How do we do a limit in a place with so many choices?!
\end{frame}

\begin{frame}{Coarse vs fine}
\begin{defn}
If $P$ and $P'$ are two partitions of $[a,b]$ with $P\subseteq P'$, we say $P$
is \textbf{coarser} than $P'$ and $P'$ is \textbf{finer} than $P$.\\
Any finer partition is called a \textbf{refinement}.
\end{defn}
\pause
\textbf{Examples:}
\begin{itemize}
\pause
\item $\{1,2,3\}$ is \textbf{finer} than $\{1,3\}$
\pause
\item $\{1, 1.1, 2, 2.5, 3\}$ is \textbf{coarser} than $\{1, 1.1, 1.5, 2, 2.5, 2.9, 3\}$
\pause
\item $\{1, 1.5, 2, 2.5, 3\}$ and $\{1, 1.2, 1.4, 1.8, 2, 2.2, 2.4, 2.8, 3\}$ are \textbf{incomparable}
\end{itemize}
\end{frame}

\begin{frame}{Riemann integrability}
\begin{defn}
We write 
$$\int_a^b f(x) dx = A$$
if for all $\epsilon > 0$, there exists a partition $P_\epsilon$ such that for any finer partition $P$
$$|R(P,f,\{t_k\}) - A| < \epsilon$$
for all choices of sample points $\{t_k\}$ of P.\\
In this case we say $f$ is \textbf{Riemann integrable} on $[a,b]$ and call the value $A$ the \textbf{Riemann integral}.
\end{defn}
\end{frame}

\begin{frame}{A first example}
\begin{prob}
Let
$$f(x) = \left\lbrace\begin{array}{cc}
1, & x < 2\\
2, & x \geq 2
\end{array}\right.$$
What is $\int_1^3 f(x) dx$?
\end{prob}
\begin{itemize}
\pause
\item Calculus experience says: $\int_1^3 f(x) dx = 3$.
\pause
\item How do we prove it? 
\end{itemize}
\textbf{INTUITION}: the "action" is near $x=2$, so we should try to make the partition denser there
\end{frame}

\begin{frame}{A first example}
Let $\epsilon > 0$. \pause(WLOG $\epsilon < 1$)\\
Choose $P_\epsilon = \{1,2-\epsilon/4,2, 3\}$\\
A refinement $P$ of $P_\epsilon$ will look like
$$P = \{x_0=1,x_1,\dots, x_{\ell}=2-\epsilon,x_{\ell+1},\dots,x_m=2,x_{m+1},\dots, x_n=3\}$$
Then for any sample points, the Riemann sum is
\begin{align*}
R(P,f,\{t_k\})
  & = \sum_{k=1}^n f(t_k)\Delta x_k\\
  & = \sum_{k=1}^\ell f(t_k)\Delta x_k + \sum_{k=\ell+1}^m f(t_k)\Delta x_k + \sum_{k=m+1}^n f(t_k)\Delta x_k\\
\end{align*}
\end{frame}

\begin{frame}{A first example}
\begin{align*}
  & = \sum_{k=1}^\ell f(t_k)\Delta x_k + \sum_{k=\ell+1}^{m-1} f(t_k)\Delta x_k + f(t_m)\Delta x_m + \sum_{k=m+1}^n f(t_k)\Delta x_k\\
  & = \sum_{k=1}^\ell 1\Delta x_k + \sum_{k=\ell+1}^{m-1} 1\Delta x_k + f(t_m)\Delta x_m + \sum_{k=m+1}^n 2\Delta x_k\\
  & = (x_\ell-x_0) + (x_{m-1}-x_{\ell}) + f(t_m)(x_m-x_{m-1}) + 2(x_n-x_{m})\\
  & = 1-\epsilon/4 + (x_{m-1}-x_{\ell}) + f(t_m)(x_m-x_{m-1}) + 2
\end{align*}
Therefore 
\begin{align*}
|R(P,f,\{t_k\})-3|
  & = |-\epsilon/4 + (x_{m-1}-x_{\ell}) + f(t_m)(x_m-x_{m-1})|\\
  & < \epsilon/4 + \epsilon/4 + 2\epsilon/4 = \epsilon
\end{align*}
\end{frame}

\begin{frame}{Riemann-Stieltjes sums}
Notation: given a real-valued function $\alpha$ on $[a,b]$\\
$$\Delta \alpha_k(P) = \alpha(x_k)-\alpha(x_{k-1}),\quad k=1,2,\dots,n.$$
\pause
\begin{defn}
\pause
Let $f$ and $\alpha$ be real-valued functions on $[a,b]$, and let $P = \{x_0,x_1,\dots,x_n\}$ be a partition of $[a,b]$.
\pause
A sum of the form
$$S(P,f,\alpha,\{t_k\}) = \sum_{k=1}^n f(t_k) \Delta\alpha_k(P),$$
for some set of sample points $\{t_k\}$ of $P$.
\end{defn}
\end{frame}

\begin{frame}{Riemann-Stieltjes integral}
\begin{defn}
\pause
Let $f$ and $\alpha$ be real-valued functions on $[a,b]$.
\pause
We say $f$ is \textbf{Riemann integrable} with respect to $\alpha$ on $[a,b]$ if there exists $A\in\mathbb{R}$ with the following property:
\pause 
for all $\epsilon > 0$, there exists a partition $P_\epsilon$ of $[a,b]$ such that for all refinements $P$ of $P_\epsilon$
$$|S(P,f,\alpha,\{t_k\}) - A| < \epsilon,$$
for all sets of sample points $\{t_k\}$ of $P$.
\end{defn}
\pause
Notation: in this case $A$ is denote by $\int_a^b f d\alpha$ or $\int_a^b f(x) d\alpha(x)$
\end{frame}

\begin{frame}{Challenge!}
\begin{prob}
Suppose that $a,b$ are real numbers with $a<b$ and let $\alpha$ be a real-valued function on $[a,b]$.\\
Prove that a constant function $f(x) = c$ is integrable with respect to $\alpha$ on $[a,b]$ and
$$\int_a^b fd\alpha = c\alpha(b)-c\alpha(a).$$
\end{prob}
\end{frame}

\begin{frame}{Challenge!}
\begin{prob}
Suppose that $a,b$ are real numbers with $a<b$ and let $\alpha$ be a constant function on $a<b$.\\
If $f(x)$ is a function on $[a,b]$, show that $f(x)$ is integrable with respect to $\alpha$ on $[a,b]$.\\
$$\text{What is $\int_a^b fd\alpha$?}$$
\end{prob}
\end{frame}

\begin{frame}{Challenge!}
\begin{prob}
Suppose that $a,b,c$ are real numbers with $a<c<b$ and let $\alpha$ be the step function
$$\alpha(x) = \left\lbrace\begin{array}{cc}
0, & x < c\\
1, & x \geq c
\end{array}\right.$$
Show that if $f(x)$ is a real function on $[a,b]$ which is continuous at $c$, then $f$ is integrable with respect to $\alpha$ on $[a,b]$.\\
$$\text{What is $\int_a^b fd\alpha$?}$$
\end{prob}
\end{frame}


\begin{frame}{Linearity of integration}
\begin{thm}
Suppose that $f$ and $g$ are Riemann integrable with respect to $\alpha$ on $[a,b]$.\\
\pause
Then for any $c_1,c_2\in\mathbb{R}$, $c_1f + c_2g$ is Riemann integrable with respect to $\alpha$ and
$$\int_a^b (c_1f + c_2g) d\alpha = c_1\int_a^b f d\alpha + c_2\int_a^b g d\alpha.$$
\end{thm}

\end{frame}
\begin{frame}{Linearity of integration}
\begin{proof}
\pause
Let $\epsilon > 0$ and WLOG assume $c_1,c_2\neq 0$.\\
\pause
Choose a partition $P_\epsilon'$ such that
$$|S(P,f,\alpha,\{t_k\}) - \int_a^bf d\alpha| < \epsilon/2|c_1|,$$
for all refinements $P$ of $P_\epsilon'$ and sets of sample points $\{t_k\}$ of $P$.\\
\pause
Also, choose a partition $P_\epsilon''$ such that
$$|S(P,g,\alpha,\{t_k\}) - \int_a^bg d\alpha| < \epsilon/2|c_2|,$$
for all refinements $P$ of $P_\epsilon''$ and sets of sample points $\{t_k\}$ of $P$.\\
\pause
Let $P_\epsilon = P_\epsilon'\cup P_\epsilon''$. 
\pause
Then any refinement $P$ of $P_\epsilon$, is a refinement of $P_\epsilon'$ and $P_\epsilon''$.
\end{proof}
\end{frame}
\begin{frame}{Linearity of integration}
\begin{proof}
Therefore if $P$ is a refinement of $P_\epsilon$, and $\{t_k\}$ is a set of sample points of $P$, then by the triangle inequality
\pause
\begin{align*}
|S(P,c_1f+c_2g,\alpha,\{t_k\}) - (c_1\int_a^b fd\alpha+c_2\int_a^bgd\alpha)|\\
 = |c_1S(P,f,\alpha,\{t_k\}) + c_2S(P,g,\alpha,\{t_k\})- (c_1\int_a^b fd\alpha+c_2\int_a^bgd\alpha)|\\
 = |c_1S(P,f,\alpha,\{t_k\}) - c_1\int_a^b fd\alpha + c_2S(P,g,\alpha,\{t_k\})- c_2\int_a^bgd\alpha|\\
 \leq |c_1|\cdot |S(P,f,\alpha,\{t_k\}) - \int_a^b fd\alpha| + |c_2|\cdot |S(P,g,\alpha,\{t_k\})- \int_a^bgd\alpha|\\
 < \epsilon/2 + \epsilon/2 = \epsilon.
\end{align*}
\pause
Since $\epsilon>0$ was arbitrary, this proves the theorem.
\end{proof}
\end{frame}

\begin{frame}{Linearity of integration}
\begin{thm}
Suppose that $f$ is Riemann integrable with respect to both $\alpha$ and $\beta$ on $[a,b]$.\\
\pause
Then for any $c_1,c_2\in\mathbb{R}$, $f$ is Riemann integrable with respect to $c_1\alpha + c_2\beta$ and
$$\int_a^b f d(c_1\alpha + c_2\beta) = c_1\int_a^b f d\alpha + c_2\int_a^b f d\beta.$$
\end{thm}
\end{frame}

\begin{frame}{Challenge!}
\begin{prob}
Prove the last theorem!
\end{prob}
\end{frame}

\begin{frame}{Linearity of integration}
\begin{thm}
Suppose that $f$ is a real-valued function on $[a,b]$ and that $c\in (a,b)$.
\pause
If $\int_a^c fd\alpha$ and $\int_b^c fd\alpha$ exist, then $\int_a^cfd\alpha$ exists and
$$\int_a^c fd\alpha + \int_c^b fd\alpha = \int_a^b fd\alpha.$$
\pause
Conversely, if $\int_a^c fd\alpha$ exists then $\int_a^c fd\alpha$ and $\int_b^c fd\alpha$ exist and
$$\int_a^c fd\alpha + \int_c^b fd\alpha = \int_a^b fd\alpha.$$
\end{thm}
\end{frame}

\begin{frame}{Integration by parts}
The integrand and integrator play dual roles:
\begin{thm}{Integration by parts}
Let $f$ be integrable with respect to $\alpha$ on $[a,b]$.\\
Then $\alpha$ is integrable with respect to $f$ on $[a,b]$ and
$$\int_a^b f d\alpha = f(b)\alpha(b)-f(a)\alpha(a)-\int_a^b \alpha df.$$
\end{thm}
\end{frame}

\begin{frame}{Relation to Riemann integrals}
\begin{thm}
Let $f$ be integrable with respect to $\alpha$ on $[a,b]$.\\
If $\alpha$ is continuous on $[a,b]$ and continuously differentiable on $(a,b)$, then
$f(x)\alpha'(x)$ is integrable with respect to $x$ on $[a,b]$ and
$$\int_a^b f d\alpha = \int_a^bf(x)\alpha'(x)dx.$$
\end{thm}
\end{frame}


\end{document}


